% 数学结论及完整证明已移入正文相应公式之后。


% 数学结论及完整证明已移入正文相应公式之后。


% 数学结论及完整证明已移入正文相应公式之后。


% 数学结论及完整证明已移入正文相应公式之后。


% 数学结论及完整证明已移入正文相应公式之后。


% 数学结论及完整证明已移入正文相应公式之后。


\begin{derivation}[从\eqnrefs{eq:pinem_finite_pulse_local_phase}到\eqnrefs{eq:park_Cjn_general}]
固定一个电子切片 $(\zeta,\tau)$，把 $\beta(\zeta,\tau)$ 暂记为 $\beta$，并令
\[
 \theta\equiv\kappa_\omega\zeta.
\]
局域散射算符是纯相位
\begin{equation*}
 U(\theta)=\exp\!\left(\beta \ee^{\ii\theta}-\beta^*\ee^{-\ii\theta}\right).
\end{equation*}
由于两个指数中的量只是复数，它们彼此对易，所以可精确分解为
\begin{equation*}
 U(\theta)
 =\ee^{\beta \ee^{\ii\theta}}
  \ee^{-\beta^*\ee^{-\ii\theta}}.
\end{equation*}
分别作 Taylor 展开，得到双重级数
\begin{align*}
 U(\theta)
 &=\left[\sum_{r=0}^{\infty}
 \frac{\beta^r}{r!}\ee^{\ii r\theta}\right]
 \left[\sum_{s=0}^{\infty}
 \frac{(-\beta^*)^s}{s!}\ee^{-\ii s\theta}\right]\\
 &=\sum_{r,s\ge0}
 \frac{\beta^r(-\beta^*)^s}{r!s!}
 \ee^{\ii(r-s)\theta}.
\end{align*}
这里 $r$ 与 $s$ 分别计数从两个指数中取出的正、负频因子数；真正由能量测量分辨的是它们的差
\[
 \ell=r-s,
\]
而不是 $r$、$s$ 各自。于是同一个最终边带 $\ell$ 包含无穷多条具有相同净交换量、但往返次数不同的振幅路径。

若 $\ell\ge0$，所有非负整数解为
\[
 r=j+\ell,\qquad s=j,
 \qquad j=0,1,2,\ldots;
\]
若 $\ell<0$，则
\[
 r=j,\qquad s=j-\ell.
\]
引入
\[
 \ell_+=\max(\ell,0),
 \qquad
 \ell_-=\max(-\ell,0),
\]
两种情况统一为
\[
 r=j+\ell_+,
 \qquad
 s=j+\ell_-.
\]
因此第 $j$ 条路径对净边带 $\ell$ 的复振幅贡献为
\begin{equation*}
 C_j^{(\ell)}
 =\frac{\beta^{j+\ell_+}(-\beta^*)^{j+\ell_-}}
 {(j+\ell_+)!(j+\ell_-)!},
\end{equation*}
恢复 $(\zeta,\tau)$ 依赖即得到正文\eqnrefs{eq:park_Cjn_general}。这里 $C_j^{(\ell)}$ 是振幅而不是概率；所有具有相同 $\ell$ 的路径必须先相干相加：
\begin{equation*}
 A_\ell(\zeta,\tau)
 =\sum_{j=0}^{\infty}C_j^{(\ell)}(\zeta,\tau).
\end{equation*}

这一求和可以显式完成。写
\[
 \beta=|\beta|\ee^{\ii\varphi_\beta}.
\]
对 $\ell\ge0$，有
\begin{align*}
 A_\ell
 &=\beta^\ell
 \sum_{j=0}^{\infty}
 \frac{(-|\beta|^2)^j}{j!(j+\ell)!}\\
 &=\ee^{\ii\ell\varphi_\beta}
 \sum_{j=0}^{\infty}
 \frac{(-1)^j|\beta|^{2j+\ell}}{j!(j+\ell)!}\\
 &=\ee^{\ii\ell\varphi_\beta}J_\ell(2|\beta|),
\end{align*}
其中最后一步使用 Bessel 函数的级数定义
\[
 J_\ell(2x)
 =\sum_{j=0}^{\infty}
 \frac{(-1)^j x^{2j+\ell}}{j!(j+\ell)!}.
\]
对 $\ell<0$ 使用 $J_{-n}(x)=(-1)^nJ_n(x)$ 得到同一统一结果
\begin{equation*}
 \boxed{
 A_\ell=\ee^{\ii\ell\varphi_\beta}J_\ell(2|\beta|).}
\end{equation*}
因此 Bessel 边带并不是额外假设，而是同一净交换通道内所有往返路径振幅的精确重求和。

还可以直接检验概率守恒。由于指数
$\beta \ee^{\ii\theta}-\beta^*\ee^{-\ii\theta}$
纯虚，任意 $\theta$ 都有 $|U(\theta)|=1$。Fourier 系数定义为
\begin{equation*}
 A_\ell=\frac{1}{2\pi}
 \int_0^{2\pi}\dd\theta\,
 U(\theta)\ee^{-\ii\ell\theta}.
\end{equation*}
Parseval 恒等式于是给
\begin{align*}
 \sum_{\ell\in\mathbb Z}|A_\ell|^2
 &=\frac{1}{2\pi}\int_0^{2\pi}\dd\theta\,|U(\theta)|^2\\
 &=1.
\end{align*}
\end{derivation}

% 数学结论及完整证明已移入正文相应公式之后。


\begin{derivation}[从\eqnrefs{eq:pinem_fast_branch_envelope_ratio}到\eqnrefs{eq:feist2015_long_interaction_condition}]
把归一化开关包络取在有限区间 $[t_i,t_f]$，满足
\[
 \widetilde w(0)=\int_{t_i}^{t_f}\dd t\,w(t)=1,
 \qquad
 \widetilde w(\Omega)=\int_{t_i}^{t_f}\dd t\,w(t)\ee^{\ii\Omega t}.
\]
正文的一阶微扰已经把慢支路与快支路之比化为
$|\widetilde w(2\omega)|$。为了看清“作用时间长”为什么能够压低它，假设 $w$ 在作用区间绝对连续。对 $\Omega\ne0$ 分部积分：
\begin{align*}
 \widetilde w(\Omega)
 &=\left.\frac{w(t)\ee^{\ii\Omega t}}{\ii\Omega}\right|_{t_i}^{t_f}
 -\frac1{\ii\Omega}\int_{t_i}^{t_f}\dd t\,
 \dot w(t)\ee^{\ii\Omega t}.
\end{align*}
取模并用三角不等式得到严格上界
\begin{align*}
 |\widetilde w(\Omega)|
 \le\frac{|w(t_f)|+|w(t_i)|
 +\int_{t_i}^{t_f}\dd t\,|\dot w(t)|}{|\Omega|}.
\end{align*}
因此在 $\Omega=2\omega$ 处，若包络端点平滑降到零，便有
\begin{equation*}
 \left|\frac{\mathcal A_{\rm fast}^{(1)}}{\mathcal A_{\rm slow}^{(1)}}\right|
 =|\widetilde w(2\omega)|
 \le\frac{1}{2\omega}
 \int_{t_i}^{t_f}\dd t\,|\dot w(t)|.
\end{equation*}
这个式子比单纯说“$T$ 很长”更精确：快支路由包络在光学频率尺度上的变化率控制；同样持续时间的两个脉冲，只要边缘平滑程度不同，反旋泄漏就可以不同。

矩形窗口不满足端点为零的平滑条件，需要直接积分。对
$w(t)=T^{-1}\mathbf 1_{[-T/2,T/2]}(t)$，
\begin{align*}
 \widetilde w(2\omega)
 &=\frac1T\int_{-T/2}^{T/2}\dd t\,\ee^{2\ii\omega t}\\
 &=\frac{\sin(\omega T)}{\omega T}
 =\sinc(\omega T).
\end{align*}
所以正文\eqnrefs{eq:feist2015_long_interaction_condition} 的一般小参数是
$\epsilon_{2\omega}=|\widetilde w(2\omega)|$；只有指定矩形开关之后，它才退化为由 $\omega T$ 单独控制的 sinc 条件。
\end{derivation}

% 数学结论及完整证明已移入正文相应公式之后。


\begin{derivation}[从\eqnrefs{eq:universal_free_electron_boson_hamiltonian}到\eqnrefs{eq:qpinem-single-mode-hamiltonian-dp68}]

一般线性电子--玻色相互作用在相互作用绘景中可写成
\[
\hat V_I(t)
=
\hbar\sum_\lambda
\left[
\hat{\mathcal G}_{\lambda,I}(t)\hat a_\lambda
+
\hat{\mathcal G}_{\lambda,I}^\dagger(t)\hat a_\lambda^\dagger
\right].
\]
这里将光模自由演化的 $\ee^{-\ii\omega_\lambda t}$ 也吸收到 $\hat{\mathcal G}_{\lambda,I}$ 中；$\hat a_\lambda$ 为固定的相互作用绘景模算符，不能再重复附加该相位。$\hat{\mathcal G}_{\lambda,I}$ 仍作用在完整正能电子连续谱上。先把电子自由本征态记为 $|p,s\rangle$，则
\[
\hat{\mathcal G}_{\lambda,I}(t)
=
\sum_{ss'}\int\dd p\,\dd p'\,
\mathcal G_{\lambda;s's}(p',p)\,
\ee^{+\ii\{[E(p')-E(p)]/\hbar-\omega_\lambda\}t}
|p',s'\rangle\langle p,s|.
\]
因此“单模 QPINEM”需要依次做三个投影，而不是一步把完整 Hamiltonian 改写成升降算符。

第一步是环境谱投影。设目标模式为 $\lambda_0$，其频率为 $\omega_0$；其余模式与目标模式的最小有效失谐记为
\[
\Delta_{\rm spec}
=
\inf_{\lambda\neq\lambda_0}
|\omega_\lambda-\omega_0|.
\]
若作用时间 $T$ 和实际耦合带宽满足
\[
\Delta_{\rm spec}T\gg1,
\qquad
\sup_{\lambda\neq\lambda_0}
\int_{t_i}^{t_f}\dd t\,
\|\hat{\mathcal G}_{\lambda,I}(t)\|
\ll1,
\]
则有限时间 Fourier 核抑制非目标模式，环境和式可约化为 $\lambda_0$。

第二步是电子通道投影。连续谱广义本征态满足 Dirac--$\delta$ 归一化，而离散通道模型需要一组真正 Kronecker 正交的归一态。先在长度 $L_z$ 的周期盒中正则化纵向连续谱，使允许动量间隔为 $2\pi\hbar/L_z$；再选择与主导动量转移 $\hbar\kappa_z$ 相容的一串正交盒归一态。记
\[
|\ell\rangle_V
\equiv
|p_\ell,s_0\rangle_V,
\qquad
p_\ell=p_0+\ell\hbar\kappa_z,
\qquad
{}_V\langle\ell|m\rangle_V=\delta_{\ell m}.
\]
若实际电子是有限宽波包，也可以用彼此近似正交、中心动量为 $p_\ell$ 的归一化通道包代替盒态；最终有效格点只要求相邻通道的重叠误差小于所保留动力学精度。记号上省略盒归一化下标 $V$，而离散化极限保持不变。投影后的电子耦合算符为
\[
\hat{\mathcal G}_{0,I}(t)
=
\sum_\ell
\mathcal G_\ell(t)
\ee^{+\ii\Delta_\ell t}
|\ell+1\rangle\langle\ell|
+
\hat R_{\rm off}(t),
\]
其中
\[
\Delta_\ell
=
\frac{E(p_{\ell+1})-E(p_\ell)}{\hbar}
-\omega_0
\]
是逐边真实失谐，$\hat R_{\rm off}$ 收集自旋翻转、非主导动量转移和其它离散通道。

第三步才是平移对称近似。若显著占据窗口 $|\ell|\le L_{\rm occ}$ 内
\[
\epsilon_{\mathcal G}
\equiv
\max_\ell
\left|
\frac{\mathcal G_\ell-\mathcal G_0}{\mathcal G_0}
\right|
\ll1,
\qquad
\epsilon_\Delta
\equiv
T\max_\ell|\Delta_\ell-\Delta_0|
\ll1,
\]
便可把共同边耦合写成
\[
\mathcal G_\ell(t)\ee^{+\ii\Delta_\ell t}
=
\mathcal G(t)
+
O\!\left(|\mathcal G_0|(\epsilon_{\mathcal G}+\epsilon_\Delta)\right),
\]
并定义双边电子平移算符
\[
\hat b^\dagger
=
\sum_{\ell\in\mathbb Z}
|\ell+1\rangle\langle\ell|,
\qquad
\hat b
=
\sum_{\ell\in\mathbb Z}
|\ell\rangle\langle\ell+1|.
\]
于是
\[
\hat{\mathcal G}_{0,I}(t)
=
\mathcal G(t)\hat b^\dagger
+\delta\hat{\mathcal G}(t).
\]

把这一结果代回单模完整 Hamiltonian，
\[
\hat V_I(t)
=
\hbar
\left[
\mathcal G(t)\hat b^\dagger\hat a
+
\mathcal G^*(t)\hat b\hat a^\dagger
\right]
+\hat R_{\rm RWA}(t)
+\hat R_{\rm tr}(t),
\]
其中 $\hat R_{\rm tr}$ 收集逐边不均匀和通道投影误差。对一般实线性场，完整相互作用还包含电子上移同时产生光子、电子下移同时湮灭光子的反旋支路；它们在相互作用绘景中带有约 $2\omega_0$ 的快速相位。这些快速支路能否舍弃，不能只由“$\omega_0$ 很大”判断。前章\eqnrefs{eq:rwa_projected_dynamical_error} 已把旋波余项的真实有限时间误差写成 Duhamel 范数；在脉冲包络缓变且反旋相位频率约为 $2\omega_0$ 的常见情形，分部积分才进一步给出\eqnrefs{eq:rwa-finite-time-error-scales-dp68} 那类充分尺度条件。只有该投影传播误差满足
\[
\epsilon_{\rm RWA}^{\rm dyn}\ll1
\]
时，反旋支路才可受控消去；$|\mathcal G|/\omega_0$ 与 $(\omega_0T)^{-1}$ 只是缓变脉冲下对这个更一般误差界的估计。

保留真实矩阵元的记号 $\mathcal G(t)$，并另定义 $\upsilon(t)=-\ii\mathcal G(t)$，即 $\mathcal G(t)=\ii\upsilon(t)$。这样 $g_{\mathrm Q}=\int\upsilon(t)\dd t$ 与经典 $\beta$ 的相位对应是显式的。代入上式，
\[
\boxed{
\hat H_Q(t)
=
\ii\hbar
\left[
\upsilon(t)\hat b^\dagger\hat a
-
\upsilon^*(t)\hat b\hat a^\dagger
\right].
}
\]
这就是正文\eqnrefs{eq:qpinem-single-mode-hamiltonian-dp68}。
\end{derivation}

% 数学结论及完整证明已移入正文相应公式之后。


% 数学结论及完整证明已移入正文相应公式之后。


% 数学结论及完整证明已移入正文相应公式之后。


% 数学结论及完整证明已移入正文相应公式之后。


% 数学结论及完整证明已移入正文相应公式之后。


% 数学结论及完整证明已移入正文相应公式之后。


% 数学结论及完整证明已移入正文相应公式之后。


% 单光子吸收禁戒已移至正文质量壳失谐之后，采用命题与独立证明。

% 数学结论及完整证明已移入正文相应公式之后。


% 数学结论及完整证明已移入正文相应公式之后。


% 数学结论及完整证明已移入正文相应公式之后。


% 数学结论及完整证明已移入正文相应公式之后。


% 数学结论及完整证明已移入正文相应公式之后。

